% \iffalse meta-comment
%
% hit-fonts.dtx 配置中西文字体、字号和字形替代
%
% \fi
%
% \section{共享字体层}
%
% \changes{v4.0a}{2026/10/06}{中西文字体集可分别设置，新增按路径加载中易字体的方案；
%   等宽字体按目标字宽校正并量化到半磅，Times 方案优先使用 Consolas；pdfLaTeX 支持中易 TrueType 字形映射}
% \changes{v4.0a}{2026/10/06}{中文粗体与斜体采用 Word 对应的合成参数，宋体强调默认换楷体；
%   几何符号、温度符号和乘号使用中文字体，破折号字体可配置，封面美术字支持轮廓替代}
% \changes{v4.0a}{2026/10/06}{字号命令的可选参数改为 Word 多倍行距；移除非通行字号
%   dachu、banxiaosi、dawu，内部固定行距调用与用户字号命令分开}
%
% \option{fontset-en} 和 \option{fontset-zh} 分别选择西文和中文字体集。
% 本模块先声明函数，待字体宏包可用后由 \cs{hit@setup@fonts} 按引擎加载。
%
%    \begin{macrocode}
%<*shared-fonts>
\str_new:N \l__hit_font_en_str
\str_new:N \l__hit_font_zh_str
\tl_new:N \g__hit_font_mono_name_tl
\tl_gset:Nn \g__hit_font_mono_name_tl { Consolas }
\fp_new:N \g__hit_font_mono_scale_fp
\fp_gset:Nn \g__hit_font_mono_scale_fp { 1 }
\cs_new_protected:Npn\hit@font@kaishu@gb:{\kaishu}
\tl_new:N \g__hit_font_mac_kaiti_tl
\tl_new:N \g__hit_font_mac_fangsong_tl
\prop_new:N \g__hit_font_mac_files_prop
\prop_const_from_keyval:Nn \c__hit_font_tex_files_prop
  {FandolSong=FandolSong-Regular.otf,FandolHei=FandolHei-Regular.otf,
   FandolKai=FandolKai-Regular.otf,FandolFang=FandolFang-Regular.otf,
   TeX~Gyre~Termes=texgyretermes-regular.otf,TeX~Gyre~Heros=texgyreheros-regular.otf,
   Latin~Modern~Mono=lmmono10-regular.otf,TeX~Gyre~Termes~Math=texgyretermes-math.otf}
\tl_new:N \l__hit_font_probe_tl
% LuaTeX 的字体名数据库可能遗漏 Apple 下载字体。按 ctex 的资源目录规则查找，
% 不缓存带系统版本散列的路径；XeTeX 使用系统字体名解析。
\cs_new_protected:Npn \__hit_font_mac_prepare:
  {
    \file_if_exist:nT{/System/Library/Fonts/Menlo.ttc}
      {\bool_gset_true:N\g__hit_font_macos_bool}
    \sys_if_engine_luatex:T
      {
        \bool_if:NT\g__hit_font_macos_bool
          {
            \lua_now:n {
              local~lfs=require('lfs');
              local~function~dirs(p)
                local~ok,it,obj=pcall(lfs.dir,p);
                local~t={};
                if~ok~then~for~d~in~it,obj~do~
                  if~not(d=='.'~or~d=='..')~then~table.insert(t,d);~end;
                end;~end;
                table.sort(t);~return~t;
              end;
              local~files={['Kaiti.ttc']='g__hit_font_mac_kaiti_tl',
                ['STFANGSO.ttf']='g__hit_font_mac_fangsong_tl'};
              for~_,root~in~ipairs({'/System/Library/AssetsV2',
                '/System/Library/AssetsV2/PreinstalledAssetsV2/InstallWithOs'})~do~
                for~_,base~in~ipairs(dirs(root))~do~
                  if~base:find('^com_apple_MobileAsset_Font')~then~
                    for~_,asset~in~ipairs(dirs(root..'/'..base))~do~
                      for~file,macro~in~pairs(files)~do~
                        local~p=root..'/'..base..'/'..asset..'/AssetData/'..file;
                        if~lfs.attributes(p,'mode')=='file'~then~
                          token.set_macro(macro,root..'/'..base..'/'..asset..'/AssetData/','global');
                        end;
                      end;
                    end;
                  end;
                end;
              end;
            }
            \tl_if_empty:NF\g__hit_font_mac_kaiti_tl
              {\prop_gput:NnV\g__hit_font_mac_files_prop{Kaiti~SC}\g__hit_font_mac_kaiti_tl}
            \tl_if_empty:NF\g__hit_font_mac_fangsong_tl
              {\prop_gput:NnV\g__hit_font_mac_files_prop{STFangsong}\g__hit_font_mac_fangsong_tl}
          }
      }
  }
\cs_new_protected:Npn \__hit_font_cjk_role:n #1
  {
    \prop_get:NnN\g__hit_font_roles_prop{#1}\l__hit_font_role_name_tl
    \tl_set:Nn\l__hit_font_role_options_tl {AutoFakeBold=\c_docxlike_font_fake_bold_tl}
    \str_case:Vn\l__hit_font_role_name_tl
      {
        {FandolSong}{\tl_set:Nn\l__hit_font_role_name_tl{FandolSong-Regular.otf}
          \tl_set:Nn\l__hit_font_role_options_tl{BoldFont=FandolSong-Bold.otf}}
        {FandolHei}{\tl_set:Nn\l__hit_font_role_name_tl{FandolHei-Regular.otf}
          \tl_set:Nn\l__hit_font_role_options_tl{BoldFont=FandolHei-Bold.otf}}
        {FandolKai}{\tl_set:Nn\l__hit_font_role_name_tl{FandolKai-Regular.otf}}
        {FandolFang}{\tl_set:Nn\l__hit_font_role_name_tl{FandolFang-Regular.otf}}
        {Songti~SC}{\tl_set:Nn\l__hit_font_role_options_tl{UprightFont=*~Regular,BoldFont=*~Bold}}
        {Heiti~SC}{\tl_set:Nn\l__hit_font_role_options_tl{UprightFont=*~Light,BoldFont=*~Medium}}
        {Kaiti~SC}{\tl_set:Nn\l__hit_font_role_options_tl{UprightFont=*~Regular,BoldFont=*~Bold}}
        {Noto~Serif~CJK~SC}{\tl_set:Nn\l__hit_font_role_options_tl{BoldFont=*~Bold}}
        {Noto~Sans~CJK~SC}{\tl_set:Nn\l__hit_font_role_options_tl{BoldFont=*~Bold}}
        {AdobeSongStd-Light}{\tl_set:Nn\l__hit_font_role_options_tl{BoldFont=AdobeHeitiStd-Regular}}
        {HYShuSongEr~S}{\tl_set:Nn\l__hit_font_role_options_tl{BoldFont=HYZhongSong~S}}
        {HYZhongHei~S}{\tl_set:Nn\l__hit_font_role_options_tl{BoldFont=HYDaHei~S}}
      }
    \prop_get:NVNT\g__hit_font_mac_files_prop\l__hit_font_role_name_tl\l_tmpa_tl
      {
        \str_if_eq:VnTF\l__hit_font_role_name_tl{Kaiti~SC}
          {
            \tl_set:Nn\l__hit_font_role_name_tl{Kaiti.ttc}
            \tl_set:Nx\l__hit_font_role_options_tl{Path={\exp_not:V\l_tmpa_tl},FontIndex=0,BoldFont=Kaiti.ttc,BoldFontIndex=3}
          }
          {
            \tl_set:Nn\l__hit_font_role_name_tl{STFANGSO.ttf}
            \tl_set:Nx\l__hit_font_role_options_tl{Path={\exp_not:V\l_tmpa_tl},AutoFakeBold=\c_docxlike_font_fake_bold_tl}
          }
      }
    \__hit_font_baseline_options:n{#1}
    \iow_log:x{hithesis~CJK~#1~=~\l__hit_font_role_name_tl~[\l__hit_font_role_options_tl]}
  }
% 测量常规字面；RawFeature 将同一偏移关联到该角色的所有字重与字号。
\cs_new_protected:Npn\__hit_font_baseline_options:n #1
  {
    \sys_if_engine_luatex:T
      {
        \prop_get:NnNT\c__hit_font_baseline_reference_prop{#1}\l__hit_font_baseline_reference_tl
          {
            \group_begin:
              \normalfont\fontsize{100bp}{120bp}\selectfont
              \use:e{\exp_not:N\fontspec{\exp_not:V\l__hit_font_role_name_tl}
                [\exp_not:V\l__hit_font_role_options_tl]}
              \exp_args:NnV\docxlike_font_baseline_probe:nn{#1}\l__hit_font_baseline_reference_tl
            \group_end:
            \tl_put_right:Nx\l__hit_font_role_options_tl
              {,RawFeature={+docxbaseline=ref\g_docxlike_baseline_id_tl}}
          }
      }
  }
\cs_new_protected:Npn \__hit_font_cjk_declare:nn #1#2
  {
    \__hit_font_cjk_role:n{#2}
    \use:e{\exp_not:N\setCJKfamilyfont{#1}{\exp_not:V\l__hit_font_role_name_tl}
      [\exp_not:V\l__hit_font_role_options_tl]}
  }
\prg_new_protected_conditional:Npnn \__hit_font_family_available:n #1 {TF}
  {
    \str_if_eq:nnTF{#1}{outline}{\prg_return_true:}
    {\prop_if_in:NnTF\g__hit_font_mac_files_prop{#1}{\prg_return_true:}
      {
        \prop_get:NnNTF\c__hit_font_tex_files_prop{#1}\l__hit_font_probe_tl
          {\exp_args:NV\fontspec_font_if_exist:nTF\l__hit_font_probe_tl{\prg_return_true:}{\prg_return_false:}}
          {
            \fontspec_font_if_exist:nTF{#1}{\prg_return_true:}
              {
                \str_if_eq:nnTF{#1}{Kaiti~SC}
                  {\fontspec_font_if_exist:nTF{Kaiti~SC~Regular}{\prg_return_true:}{\prg_return_false:}}
                  {\file_if_exist:nTF{#1-Regular.otf}{\prg_return_true:}{\prg_return_false:}}
              }
          }
      }}
  }
\cs_new_protected:Npn \__hit_font_code_cjk_select:
  {
    \tl_gclear:N\g__hit_font_code_cjk_name_tl
    \clist_map_inline:Nn\c__hit_font_code_cjk_roles_clist
      {
        \prop_get:NnN\g__hit_font_roles_prop{##1}\l__hit_font_role_name_tl
        \exp_args:NV\__hit_font_family_available:nTF\l__hit_font_role_name_tl
          {
            \tl_gset_eq:NN\g__hit_font_code_cjk_name_tl\l__hit_font_role_name_tl
            \__hit_font_cjk_role:n{##1}
            \clist_map_break:
          }{}
      }
    \tl_if_empty:NT\g__hit_font_code_cjk_name_tl
      {\__hit_font_cjk_role:n{fangsong}\tl_gset_eq:NN\g__hit_font_code_cjk_name_tl\l__hit_font_role_name_tl}
    \prop_get:NnN\g__hit_font_roles_prop{fangsong}\l_tmpa_tl
    \tl_if_eq:NNF\l_tmpa_tl\g__hit_font_code_cjk_name_tl
      {\msg_warning:nnxx{hithesis}{font-code-fallback}{\l_tmpa_tl}{\g__hit_font_code_cjk_name_tl}}
  }
\cs_new_protected:Npn \__hit_set_font_zh_roles:
  {
    \__hit_font_cjk_role:n{songti}
    \use:e{\exp_not:N\setCJKmainfont{\exp_not:V\l__hit_font_role_name_tl}[\exp_not:V\l__hit_font_role_options_tl]}
    \__hit_font_cjk_role:n{heiti}
    \use:e{\exp_not:N\setCJKsansfont{\exp_not:V\l__hit_font_role_name_tl}[\exp_not:V\l__hit_font_role_options_tl]}
    \__hit_font_code_cjk_select:
    \use:e{\exp_not:N\setCJKmonofont{\exp_not:V\l__hit_font_role_name_tl}[\exp_not:V\l__hit_font_role_options_tl]}
    \__hit_font_cjk_declare:nn{zhsong}{songti}
    \__hit_font_cjk_declare:nn{zhhei}{heiti}
    \__hit_font_cjk_declare:nn{zhkai}{kaishu}
    \__hit_font_cjk_declare:nn{zhkaigb}{kaishu-gb2312}
    \cs_set_protected:Npn\hit@font@kaishu@gb:{\CJKfamily{zhkaigb}}
    \__hit_font_cjk_declare:nn{zhfs}{fangsong}
  }
\cs_new_protected:Npn \__hit_set_font_en_roles:
  {
    \prop_get:NnN\g__hit_font_roles_prop{serif}\l__hit_font_role_name_tl
    \exp_args:NNV\__hit_font_latin_set:Nn\setmainfont\l__hit_font_role_name_tl
    \prop_get:NnN\g__hit_font_roles_prop{sans}\l__hit_font_role_name_tl
    \exp_args:NNV\__hit_font_latin_set:Nn\setsansfont\l__hit_font_role_name_tl
    \prop_get:NnN\g__hit_font_roles_prop{mono}\g__hit_font_mono_name_tl
    \exp_args:NNV\__hit_font_latin_set:Nn\setmonofont\g__hit_font_mono_name_tl
    \fp_gset:Nn\g__hit_font_mono_scale_fp{1}
  }
% 最终回退按 TeX Live 文件名加载，不依赖操作系统字体注册。
\cs_new_protected:Npn \__hit_font_latin_set:Nn #1#2
  {
    \str_case:nnF{#2}
      {
        {TeX~Gyre~Termes}{#1{texgyretermes}[Extension=.otf,UprightFont=*-regular,
          BoldFont=*-bold,ItalicFont=*-italic,BoldItalicFont=*-bolditalic]}
        {TeX~Gyre~Heros}{#1{texgyreheros}[Extension=.otf,UprightFont=*-regular,
          BoldFont=*-bold,ItalicFont=*-italic,BoldItalicFont=*-bolditalic]}
        {Latin~Modern~Mono}{#1{lmmono10}[Extension=.otf,UprightFont=*-regular,
          ItalicFont=*-italic,BoldFont=lmmonolt10-bold,BoldItalicFont=lmmonolt10-boldoblique]}
      }
      {#1{#2}}
  }
\cs_new_protected:Npn \hit@setup@fonts
  {
    \exp_args:Ne\__hit_font_baseline_enable:n{\use:c{hit@font@baseline}}
    \sys_if_engine_pdftex:TF
      { \__hit_setup_fonts_pdftex: }
      { \__hit_setup_fonts_unicode: }
    \__hit_setup_logo_font:
  }
% 标识使用自带小型大写字形的字体，避免当前正文缺少 sc 字形时静默回退。
\cs_new_protected:Npn \__hit_setup_logo_font:
  {
    \sys_if_engine_pdftex:TF
      {
        \cs_new_protected:Npn \hit@logo@font
          { \fontencoding{T1}\fontfamily{qtm}\selectfont }
      }
      {
        \newfontfamily \hit@logo@font { texgyretermes }
          [ Extension = .otf ,
            UprightFont = *-regular , BoldFont = *-bold ,
            ItalicFont = *-italic , BoldItalicFont = *-bolditalic ]
      }
  }
\cs_new_protected:Npn \__hit_setup_fonts_unicode:
  {
    \__hit_font_mac_prepare:
    \__hit_font_preset_resolve:
%    \end{macrocode}
%
% 西文分项的 auto 继承整套预设；旧 times、termes 名称转换为对应角色。
%    \begin{macrocode}
\str_set:Nx \l__hit_font_en_str { \hit@font@en }
\str_if_eq:VnT \l__hit_font_en_str { } { \str_set:Nn \l__hit_font_en_str { auto } }
\str_if_eq:VnT \l__hit_font_en_str { auto }
  {
    \prop_if_in:NVTF\g__hit_font_presets_prop\g__hit_font_preset_tl
      {\str_set:Nn\l__hit_font_en_str{roles}}
      {\str_if_eq:VnTF\g__hit_font_preset_tl{auto}
        {\str_set:Nn\l__hit_font_en_str{roles}}
        {\str_set:Nn\l__hit_font_en_str{times}}}
  }
\__hit_font_canonical:N\l__hit_font_en_str
\prop_if_in:NVT\g__hit_font_presets_prop\l__hit_font_en_str
  {
    \__hit_font_preset_section:Vn\l__hit_font_en_str{serif,sans,mono}
    \str_set:Nn\l__hit_font_en_str{roles}
  }
\str_if_eq:VnT \l__hit_font_en_str { times }
  {
    \fontspec_font_if_exist:nF { Times~New~Roman }
      { \str_set:Nn \l__hit_font_en_str { termes } }
  }
\str_case:Vn\l__hit_font_en_str
  {
    {times}{
      \prop_gput:Nnn\g__hit_font_roles_prop{serif}{Times~New~Roman}
      \prop_gput:Nnn\g__hit_font_roles_prop{sans}{Arial}
      \prop_gput:Nnn\g__hit_font_roles_prop{mono}{Consolas}
      \str_set:Nn\l__hit_font_en_str{roles}}
    {termes}{
      \prop_gput:Nnn\g__hit_font_roles_prop{serif}{TeX~Gyre~Termes}
      \prop_gput:Nnn\g__hit_font_roles_prop{sans}{TeX~Gyre~Heros}
      \prop_gput:Nnn\g__hit_font_roles_prop{mono}{TeX~Gyre~Cursor}
      \str_set:Nn\l__hit_font_en_str{roles}}
  }
\clist_map_inline:nn{serif,sans,mono}
  {\prop_get:NnNT\g__hit_font_overrides_prop{##1}\l_tmpa_tl{\prop_gput:NnV\g__hit_font_roles_prop{##1}\l_tmpa_tl}}
\group_begin:
  \bool_set_false:N\l__hit_font_auto_bool
  \__hit_font_roles_fallback:n{serif,sans,mono}
\group_end:
%    \end{macrocode}
%
% 等宽字体在正文样式就绪后统一校正；none 不改动用户字体。
%    \begin{macrocode}
\cs_new_protected:Npn \__hit_set_font_en_none: { }
%    \end{macrocode}
%
% 取值无效时报错。
% 顺手把“再设一遍西文字体”这件事封成 \cs{hit@fonts@redo@en:}。
% 西文网格字距通过 fontspec 字体特性实现，需要重新声明字体。
% 在 docxlike/latin-font-reload 钩子中复用已选字体方案，使正文和重载后的字体设置一致。
%    \begin{macrocode}
\cs_if_exist:cTF { __hit_set_font_en_ \l__hit_font_en_str : }
  {
    \use:c { __hit_set_font_en_ \l__hit_font_en_str : }
    \cs_new_protected:Npn \hit@fonts@redo@en:
      { \use:c { __hit_set_font_en_ \l__hit_font_en_str : } }
    \hook_gput_code:nnn { docxlike / latin-font-reload } { hithesis }
      { \hit@fonts@redo@en: }
  }
  { \msg_error:nnx { hithesis } { font-en-unknown } { \l__hit_font_en_str } }
%    \end{macrocode}
%
% 两档按\emph{文件路径}加载中易字体，给“字体文件在、但 fontconfig 没登记”
% 的场合用：\texttt{windows-local} 读 \option{windows-font-dir} 指的目录
% （默认当前目录，用户把 \file{Simsun.ttc} 之类拷在文稿旁边就能用），
% \texttt{mac-word} 读 macOS 上 Microsoft Word 自带的那个字体目录。
%
% 文件名两套都试：新一点的 Windows 是 \file{Kaiti.ttf}、\file{Fangsong.ttf}，
% 老的是 \file{Simkai.ttf}、\file{Simfang.ttf}。
%
% 黑体是 \file{SimHei.ttf}，不是 \file{.ttc}：Word 自带的那个目录里是它，
% \pkg{ctex} 的 \texttt{windows} 档写的也是 \file{simhei.ttf}。
%
% 这一段的结构与文件名取自 \pkg{thuthesis}（LPPL，同许可证）。
%
% \emph{伪粗体系数。} 宋体没有粗体，粗宋走 \pkg{xeCJK} 的伪粗体，系数取 \pkg{docxlike} 的
% \cs{c\_docxlike\_font\_fake\_bold\_tl}（2.85，Word 描边合成的线宽，量法在那边）；上游的 2 没有对过 Word。
% 楷体、仿宋也没有粗体字面，一样开（深圳、威海封面的校名是粗楷体）。\hologo{LuaLaTeX} 上字形本来由
% \pkg{docxlike} 换成描边的，不开的话 NFSS 取不到 \texttt{b} 那一形，每编一次报一条 Font shape undefined。
%
% 粗宋体强制使用常规宋体描边，避免 BoldFont 映射到黑体；西文仍使用真实粗体。
%    \begin{macrocode}
\cs_new_protected:Npn \hit@font@fakebold@zh:
  {
    \addCJKfontfeatures
      { BoldFont = { } , AutoFakeBold = \c_docxlike_font_fake_bold_tl }
    \bfseries
  }
\tl_const:Nn \c__hit_font_mac_word_dir_tl
  { /Applications/Microsoft~Word.app/Contents/Resources/DFonts }
\cs_new_protected:Npn \__hit_set_font_zh_by_dir:n ##1
  {
    \file_if_exist:nTF { ##1 /Kaiti.ttf }
      {
        \setCJKmainfont { SimSun }
          [ Path = ##1 / , Extension = .ttc , AutoFakeBold = \c_docxlike_font_fake_bold_tl ]
        \setCJKsansfont { SimHei }
          [ Path = ##1 / , Extension = .ttf , AutoFakeBold = \c_docxlike_font_fake_bold_tl ]
        \setCJKmonofont { Fangsong }
          [ Path = ##1 / , Extension = .ttf , AutoFakeBold = \c_docxlike_font_fake_bold_tl ]
        \setCJKfamilyfont { zhsong } { SimSun }
          [ Path = ##1 / , Extension = .ttc , AutoFakeBold = \c_docxlike_font_fake_bold_tl ]
        \setCJKfamilyfont { zhhei } { SimHei }
          [ Path = ##1 / , Extension = .ttf , AutoFakeBold = \c_docxlike_font_fake_bold_tl ]
        \setCJKfamilyfont { zhkai } { Kaiti }
          [ Path = ##1 / , Extension = .ttf , AutoFakeBold = \c_docxlike_font_fake_bold_tl ]
        \setCJKfamilyfont { zhfs } { Fangsong }
          [ Path = ##1 / , Extension = .ttf , AutoFakeBold = \c_docxlike_font_fake_bold_tl ]
      }
      {
        \setCJKmainfont { SimSun }
          [ Path = ##1 / , Extension = .ttc , AutoFakeBold = \c_docxlike_font_fake_bold_tl ]
        \setCJKsansfont { SimHei }
          [ Path = ##1 / , Extension = .ttf , AutoFakeBold = \c_docxlike_font_fake_bold_tl ]
        \setCJKmonofont { Simfang }
          [ Path = ##1 / , Extension = .ttf , AutoFakeBold = \c_docxlike_font_fake_bold_tl ]
        \setCJKfamilyfont { zhsong } { SimSun }
          [ Path = ##1 / , Extension = .ttc , AutoFakeBold = \c_docxlike_font_fake_bold_tl ]
        \setCJKfamilyfont { zhhei } { SimHei }
          [ Path = ##1 / , Extension = .ttf , AutoFakeBold = \c_docxlike_font_fake_bold_tl ]
        \setCJKfamilyfont { zhkai } { Simkai }
          [ Path = ##1 / , Extension = .ttf , AutoFakeBold = \c_docxlike_font_fake_bold_tl ]
        \setCJKfamilyfont { zhfs } { Simfang }
          [ Path = ##1 / , Extension = .ttf , AutoFakeBold = \c_docxlike_font_fake_bold_tl ]
      }
  }
\cs_new_protected:cpn { __hit_set_font_zh_windows-local: }
  { \exp_args:Nx \__hit_set_font_zh_by_dir:n { \hit@windows@font@dir } }
\cs_new_protected:cpn { __hit_set_font_zh_mac-word: }
  { \exp_args:NV \__hit_set_font_zh_by_dir:n \c__hit_font_mac_word_dir_tl }
%    \end{macrocode}
%
%    \begin{macrocode}
\cs_new_protected:Npn \__hit_set_font_zh_none: { }
%    \end{macrocode}
%
% 中文字体由角色表设置；分项预设覆盖整套预设，显式角色覆盖优先。
%
% 中文斜体使用 docxlike 的合成倾斜参数，西文使用字体自带的斜体字面。
%
% 中文的斜体一律斜切，宋体也不例外；宋体里的 \cs{emph} 另换楷体正体，由
% \option{emph-songti} 管，具体由 \cs{hit@font@italic@pass:} 在正文开始前应用。
%
% 倾斜特性通过 defaultCJKfontfeatures 应用于各中文字族。
%
%    \begin{macrocode}
\str_set:Nx \l__hit_font_zh_str { \hit@font@zh }
\str_if_eq:VnT\l__hit_font_zh_str{}{\str_set:Nn\l__hit_font_zh_str{auto}}
\str_if_eq:VnF \l__hit_font_zh_str { }
  {
    \str_if_eq:VnT \l__hit_font_zh_str { auto }
      {\prop_if_in:NVTF\g__hit_font_presets_prop\g__hit_font_preset_tl
        {\str_set:Nn\l__hit_font_zh_str{roles}}
        {\str_set:Nx\l__hit_font_zh_str{\g__hit_font_preset_tl}}}
    \__hit_font_canonical:N\l__hit_font_zh_str
    \str_if_eq:VnT\l__hit_font_zh_str{auto}{\str_set:Nn\l__hit_font_zh_str{roles}}
    \prop_if_in:NVT\g__hit_font_presets_prop\l__hit_font_zh_str
      {
        \__hit_font_preset_section:Vn\l__hit_font_zh_str{songti,heiti,kaishu,fangsong,kaishu-gb2312,lishu,xinwei}
        \str_set:Nn\l__hit_font_zh_str{roles}
      }
    \prop_map_inline:Nn\g__hit_font_roles_prop
      {\iow_log:x{hithesis~font~role~##1~=~##2}}
    \defaultCJKfontfeatures { AutoFakeSlant = \c_docxlike_font_fake_slant_tl }
    \cs_if_exist:cTF { __hit_set_font_zh_ \l__hit_font_zh_str : }
      { \use:c { __hit_set_font_zh_ \l__hit_font_zh_str : } }
      { \msg_error:nnx { hithesis } { font-zh-unknown } { \l__hit_font_zh_str } }
%    \end{macrocode}
%
% 中文字族命令由模板统一提供。
%    \begin{macrocode}
    \cs_set_protected:Npn \songti { \CJKfamily { zhsong } }
    \cs_set_protected:Npn \heiti { \CJKfamily { zhhei } }
    \cs_set_protected:Npn \fangsong { \CJKfamily { zhfs } }
    \cs_set_protected:Npn \kaishu { \CJKfamily { zhkai } }
  }% \str_if_eq:VnF \l__hit_font_zh_str { } 到此为止
%    \end{macrocode}
%
% 哪些字按东亚提示排见下面的 \cs{\_\_hit\_font\_hints:}，\hologo{pdfLaTeX} 那一路也调它。
%    \begin{macrocode}
    \__hit_font_hints:
%    \end{macrocode}
%
% 斜体那一遍重声明也挂在 \texttt{begindocument} 上：\option{emph-songti}
% 在导言区随时能写，字族也要等 \pkg{ctex} 与模板都声明完才齐。
%    \begin{macrocode}
    \AddToHook { begindocument } { \hit@font@italic@pass: }
  }% \hit@setup@fonts 到此为止
%    \end{macrocode}
%
% \begin{macro}{\__hit_font_hints:}
% 按东亚提示排的字，两条字体路径（\cs{\_\_hit\_setup\_fonts\_unicode:} 与 \hologo{pdfLaTeX} 那一路）都调。
% 做法在 \pkg{docxlike}（Word 的 \texttt{w:hint}），这里只说哪些字。
%
% 几何符号使用中文字体，避免西文字体缺少占位方框等字形。
%    \begin{macrocode}
\cs_new_protected:Npn \__hit_font_hints:
  {
    \docxlike_font_hint_east_asian:n { "25A1 , "25A0 , "25CB , "25CF , "25B3 , "25B2 }
%    \end{macrocode}
% ℃（U+2103）与 ℉（U+2109）同理：Times New Roman 里没有这两个字，Unicode 又把
% 它们归在“字母式符号”区块、算西文，范例表题“温度：T=16 ℃”里的 ℃ 就排丢了。
% 范例的 docx 里这个字带 \texttt{w:hint="eastAsia"}，Word 用宋体排、占一格。
%    \begin{macrocode}
    \docxlike_font_hint_east_asian:n { "2103 , "2109 }
%    \end{macrocode}
%
% 破折号归哪一副字看 \option{em-dash}，键的依据见
% \file{src/engine/hit-keylist.dtx}。
%
% 乘号 \texttt{×}（U+00D7）一律按汉字排，取宋体、占一格。Unicode 把它归在拉丁补充区，
% \pkg{xeCJK} 与 \pkg{luatexja} 默认都交给西文字体；Word 里中文输入法敲出来的 \texttt{×}
% 带 \texttt{w:hint="eastAsia"}，排的是宋体，范例分类号的「××××」、致谢与简历里的占位叉
% 都是这样（2026-09-27 定为默认）。公式里的 \cs{times} 不受影响，走数学字体。
% \pkg{docxlike} 在 \texttt{begindocument} 上把它声明进半角标点类，所以同样发在
% \texttt{begindocument/end} 上。
%
% 字符归类必须晚于 docxlike 的标点初始化，因此使用 begindocument/end。
%    \begin{macrocode}
    \AddToHook { begindocument / end }
      {
        \bool_if:NT \g__hit_em_dash_latin_bool
          { \docxlike_font_hint_default:n { "2014 } }
        \docxlike_font_hint_east_asian:n { "D7 }
      }
  }
%    \end{macrocode}
%
% \end{macro}
%
%
% \subsection{\hologo{pdfLaTeX} 的字体}
%
% \hologo{pdfTeX} 不认 OpenType，\pkg{fontspec}、\pkg{xeCJK}、\pkg{luatexja} 都用不了。中文走
% \pkg{ctex} 的 \hologo{pdfTeX} 后端（\pkg{CJKutf8}），字形由 \pkg{zhmetrics} 的
% \texttt{unisong}、\texttt{unihei}、\texttt{unikai}、\texttt{unifs} 几套度量映射到中易的
% TrueType 字体文件上，用的与 \hologo{XeLaTeX}、\hologo{LuaLaTeX} 是同一批文件。做法参考
% \pkg{gxuthesis}：
% \begin{itemize}
% \item 字体文件逐个目录去找：\option{windows-font-dir}（默认当前目录）、
%   \texttt{\$HOME/Library/Fonts}、\texttt{/Library/Fonts}、Windows 与常见 Linux 的字体目录，最后问
%   \texttt{kpsewhich}（在受限的 shell escape 白名单里）。文件名大小写几种写法都试。
% \item 路径带空格的目录不找。映射行以空格分段，带空格的路径写不进去，只能在文档结尾嵌字体时
%   才报“打不开字体文件”；Word 自带的那个目录在 \file{Microsoft Word.app} 里，所以也用不上。
% \item \pkg{ctex} 自己那份 Windows 字体映射写的是裸文件名，靠 kpathsea 找，多半找不到；
%   这里用找到的绝对路径各写一条 \cs{pdfmapline} 压过去，挂在 \texttt{shipout/firstpage} 上，
%   与 \pkg{ctex} 同一个钩子、排在它后面。
% \item 宋体一副都找不到时整个退到 \TeX~Live 自带的文鼎（\texttt{gbsn}／\texttt{gkai}），
%   没有真黑体与仿宋，外观降级但编得过；别的字体缺了退到宋体。
% \item 汉字的粗体与斜体由 \pkg{docxlike} 在字形这一层画（Word 的合成粗体与斜体），见
%   \file{src/docxlike/docxlike-format.dtx} 的“\hologo{pdfTeX} 下的汉字”。
% \end{itemize}
% 西文用 \pkg{newtxtext}（Times 的克隆），T1 编码；数学在 \file{src/engine/hit-math.dtx} 里
% 走 \pkg{newtxmath} 那一档。
%    \begin{macrocode}
\seq_new:N \l__hit_pdftex_dir_seq
\tl_new:N \l__hit_pdftex_file_tl
\tl_new:N \l__hit_pdftex_home_tl
\tl_new:N \l__hit_pdftex_song_tl
\tl_new:N \l__hit_pdftex_hei_tl
\tl_new:N \l__hit_pdftex_kai_tl
\tl_new:N \l__hit_pdftex_fs_tl
\tl_new:N \l__hit_pdftex_li_tl
\bool_new:N \g__hit_pdftex_arphic_bool
\cs_new_protected:Npn \__hit_pdftex_dirs:
  {
    \seq_clear:N \l__hit_pdftex_dir_seq
    \exp_args:Nx \tl_if_in:nnTF { \hit@windows@font@dir } { ~ }
      { \msg_warning:nnx { hithesis } { pdftex-font-dir-space } { \hit@windows@font@dir } }
      { \seq_put_right:Nx \l__hit_pdftex_dir_seq { \hit@windows@font@dir } }
    \sys_if_shell:T
      {
        \sys_get_shell:nnN { kpsewhich~-var-value=HOME } { } \l__hit_pdftex_home_tl
        \tl_trim_spaces:N \l__hit_pdftex_home_tl
        \tl_if_empty:NF \l__hit_pdftex_home_tl
          {
            \tl_if_in:NnF \l__hit_pdftex_home_tl { ~ }
              { \seq_put_right:Nx \l__hit_pdftex_dir_seq { \l__hit_pdftex_home_tl /Library/Fonts } }
          }
      }
    \seq_put_right:Nn \l__hit_pdftex_dir_seq { /Library/Fonts }
    \seq_put_right:Nn \l__hit_pdftex_dir_seq { C:/Windows/Fonts }
    \seq_put_right:Nn \l__hit_pdftex_dir_seq { /usr/share/fonts/truetype/windows }
    \seq_put_right:Nn \l__hit_pdftex_dir_seq { /usr/share/fonts/windows }
    \seq_put_right:Nn \l__hit_pdftex_dir_seq { /usr/local/share/fonts }
  }
% \__hit_pdftex_find:nN {候选文件名 clist} \tl ：逐目录逐名字找，找到存绝对路径，找不到清空
\cs_new_protected:Npn \__hit_pdftex_find:nN #1#2
  {
    \tl_clear:N #2
    \seq_map_inline:Nn \l__hit_pdftex_dir_seq
      {
        \clist_map_inline:nn {#1}
          {
            \file_if_exist:nT { ##1 / ####1 }
              {
                \tl_set:Nn #2 { ##1 / ####1 }
                \clist_map_break:n { \seq_map_break: }
              }
          }
      }
    \tl_if_empty:NT #2
      {
        \sys_if_shell:T
          {
            \clist_map_inline:nn {#1}
              {
                \sys_get_shell:nnN { kpsewhich~##1 } { } \l__hit_pdftex_file_tl
                \tl_trim_spaces:N \l__hit_pdftex_file_tl
                \tl_if_empty:NF \l__hit_pdftex_file_tl
                  {
                    \tl_set_eq:NN #2 \l__hit_pdftex_file_tl
                    \clist_map_break:
                  }
              }
          }
      }
  }
\cs_new_protected:Npn \__hit_pdftex_map:nV #1#2
  {
    \use:x
      {
        \exp_not:N \AddToHook { shipout / firstpage }
          {
            \exp_not:N \pdfmapline { = #1 @Unicode@ ~ < #2 }
            \exp_not:N \pdfmapline { = #1 sl @Unicode@ ~ < #2 }
          }
      }
  }
\cs_new_protected:Npn \__hit_setup_fonts_pdftex:
  {
    \__hit_pdftex_dirs:
    \__hit_pdftex_find:nN { simsun.ttc , Simsun.ttc , SIMSUN.TTC , simsun.ttf } \l__hit_pdftex_song_tl
    \tl_if_empty:NTF \l__hit_pdftex_song_tl
      {
        \bool_gset_true:N \g__hit_pdftex_arphic_bool
        \msg_warning:nn { hithesis } { pdftex-arphic }
        \tl_set:Nn \CJKrmdefault { gbsn }
        \tl_set:Nn \CJKsfdefault { gkai }
        \tl_set:Nn \CJKttdefault { gkai }
      }
      {
        \__hit_pdftex_find:nN { simhei.ttf , SimHei.ttf , SIMHEI.TTF } \l__hit_pdftex_hei_tl
        \__hit_pdftex_find:nN { simkai.ttf , Kaiti.ttf , KaiTi.ttf , SIMKAI.TTF } \l__hit_pdftex_kai_tl
        \__hit_pdftex_find:nN { simfang.ttf , Fangsong.ttf , FangSong.ttf , SIMFANG.TTF } \l__hit_pdftex_fs_tl
        \tl_if_empty:NT \l__hit_pdftex_hei_tl { \tl_set_eq:NN \l__hit_pdftex_hei_tl \l__hit_pdftex_song_tl }
        \tl_if_empty:NT \l__hit_pdftex_kai_tl { \tl_set_eq:NN \l__hit_pdftex_kai_tl \l__hit_pdftex_song_tl }
        \tl_if_empty:NT \l__hit_pdftex_fs_tl  { \tl_set_eq:NN \l__hit_pdftex_fs_tl  \l__hit_pdftex_kai_tl }
        \ctex_load_zhmap:nnnn { zhsong } { zhhei } { zhfs } { windows }
        \__hit_pdftex_map:nV { unisong } \l__hit_pdftex_song_tl
        \__hit_pdftex_map:nV { unihei } \l__hit_pdftex_hei_tl
        \__hit_pdftex_map:nV { unikai } \l__hit_pdftex_kai_tl
        \__hit_pdftex_map:nV { unifs } \l__hit_pdftex_fs_tl
      }
    \bool_if:NTF \g__hit_pdftex_arphic_bool
      {
        \cs_set_protected:Npn \songti { \CJKfamily { gbsn } }
        \cs_set_protected:Npn \heiti { \CJKfamily { gbsn } }
        \cs_set_protected:Npn \fangsong { \CJKfamily { gkai } }
        \cs_set_protected:Npn \kaishu { \CJKfamily { gkai } }
      }
      {
        \cs_set_protected:Npn \songti { \CJKfamily { zhsong } }
        \cs_set_protected:Npn \heiti { \CJKfamily { zhhei } }
        \cs_set_protected:Npn \fangsong { \CJKfamily { zhfs } }
        \cs_set_protected:Npn \kaishu { \CJKfamily { zhkai } }
      }
    \RequirePackage [ T1 ] { fontenc }
    \RequirePackage { newtxtext }
    \cs_gset_protected:Npn \hit@font@fakebold@zh: { \bfseries }
    \__hit_pdftex_emph_family:
    \__hit_font_hints:
  }
%    \end{macrocode}
%
% \cs{emph} 在宋体里换楷体正体（\option{emph-songti}）。\hologo{pdfTeX} 下只有
% 找到的那副楷体，写了别的字体名也用它；\texttt{fake-slant} 不声明，\cs{emph} 与
% \cs{textit} 一样斜切。字族 \texttt{hitemphsong} 的正、斜、倾斜三种形状都指向楷体的正体，标点间距照
% \texttt{zhkai} 那一族；汉字的斜切由 \pkg{docxlike} 在字形这一层做，这一族登记成真的斜体字面，不斜。
%    \begin{macrocode}
\cs_new_protected:Npn \__hit_pdftex_emph_family:
  {
    \str_if_eq:VnF \g__hit_font_emph_songti_tl { fake-slant }
      {
        \str_if_eq:VnF \g__hit_font_emph_songti_tl { kaiti }
          { \msg_warning:nn { hithesis } { pdftex-emph-font } }
        \tl_set:Nx \l__hit_pdftex_file_tl
          { \bool_if:NTF \g__hit_pdftex_arphic_bool { gkaiu } { unikai } }
        \DeclareFontFamily { C70 } { hitemphsong } { \tex_hyphenchar:D \tex_font:D = -1 ~ }
        \clist_map_inline:nn { n , it , sl }
          {
            \use:x
              {
                \exp_not:N \DeclareFontShape { C70 } { hitemphsong } { m } {##1}
                  { <-> CJK * \l__hit_pdftex_file_tl } { \exp_not:N \CJKnormal }
                \exp_not:N \DeclareFontShape { C70 } { hitemphsong } { b } {##1}
                  { <-> CJKb * \l__hit_pdftex_file_tl } { \exp_not:N \CJKbold }
                \exp_not:N \DeclareFontShape { C70 } { hitemphsong } { bx } {##1}
                  { <-> CJKb * \l__hit_pdftex_file_tl } { \exp_not:N \CJKbold }
              }
          }
        \ctex_punct_map_family:nn { hitemphsong } { zhkai }
        \docxlike_font_real_italic:n { hitemphsong }
        \bool_gset_true:N \g__hit_font_emph_bool
      }
  }
%    \end{macrocode}
%
% \begin{macro}{\hit@font@italic@pass:}
% 重声明字族以更新斜体特性；ctex 的 LuaTeX 后端须先清除字族缓存。
% \begin{itemize}
% \item 宋体（正文字族与 \texttt{zhsong}，也就是 \cs{songti} 选的那一族）：把
%   \texttt{ItalicFont}、\texttt{AutoFakeSlant} 一类全清掉，写 \texttt{AutoFakeSlant}。
%   \pkg{ctex} 的 \texttt{windows} 等字体集给宋体写了 \texttt{ItalicFont=KaiTi}，
%   不清的话 \cs{textit} 下仍是楷体。
% \item 别的字族：没有 \texttt{ItalicFont} 也没有 \texttt{AutoFakeSlant} 的补一句
%   \texttt{AutoFakeSlant}，写了的不动——用户自己 \cs{setCJKfamilyfont} 指定的
%   斜体面照旧。
% \end{itemize}
% \texttt{Script=CJK} 那一项声明时会自己加，表里取出来的那份要去掉，不然重声明后是两份。
% 特性表没变的字族不重声明，字体实例一个不多。子字族（名字带斜杠的 \texttt{FallBack} 那种）
% 跳过，主字族声明时会带上。重声明本是这里有意做的，“Redefining CJKfamily”那句警告在这一遍
% 里关掉。
%
% 然后按 \option{emph-songti} 声明 \cs{emph} 用的那一族 \texttt{hitemphsong}：
% 默认取 \texttt{zhkai} 那一族的字体与特性，写了字体名就用那副字；两种都把 \texttt{ItalicFont}
% 指回同一副字，于是 \cs{emph} 里西文照常变斜、中文是正体；\texttt{AutoFakeSlant} 另写成零，
% \pkg{luatexja} 那边默认特性里的斜切会连 \texttt{ItalicFont} 也斜掉。写 \texttt{fake-slant} 或者没有
% 楷体那一族时不声明，\cs{emph} 与 \cs{textit} 一样。
%
% \emph{楷体的扩展名要单独看。} \texttt{windows-local} 与 \texttt{mac-word} 两档按
% 路径加载，宋体是 \file{Simsun.ttc}、楷体是 \file{Kaiti.ttf}，\texttt{Extension}
% 不同。\texttt{hitemphsong} 自己带着楷体那族的 \texttt{Extension}，\texttt{ItalicFont}
% 跟着它取，写字体名本身就行。
%    \begin{macrocode}
\seq_new:N \l__hit_font_families_seq
\tl_new:N \l__hit_font_name_tl
\tl_new:N \l__hit_font_emph_name_tl
\clist_new:N \l__hit_font_emph_options_clist
\clist_new:N \l__hit_font_options_clist
\clist_new:N \l__hit_font_new_options_clist
\bool_new:N \l__hit_font_has_italic_bool
\bool_new:N \g__hit_font_emph_bool
\cs_new_protected:Npn \hit@font@italic@pass:
  {
    \cs_if_exist:NTF \g__xeCJK_family_font_name_prop
      {
        \__hit_font_italic_pass:NNNnn
          \g__xeCJK_family_font_name_prop \g__xeCJK_family_font_options_prop
          \xeCJK_set_family:nnn { xeCJK } { CJKfamily-redef }
      }
      {
        \sys_if_engine_luatex:T
          {
            \__hit_font_italic_pass:NNNnn
              \g__hit_cjk_family_font_name_prop \g__hit_cjk_family_font_options_prop
              \__hit_cjk_set_family:nnn { hithesis } { cjk-family-unknown }
          }
      }
  }
\cs_new_protected:Npn \__hit_font_italic_pass:NNNnn #1#2#3#4#5
  {
    \seq_clear:N \l__hit_font_families_seq
    \prop_map_inline:Nn #1
      { \seq_put_right:Nn \l__hit_font_families_seq {##1} }
    \group_begin:
    \msg_redirect_name:nnn {#4} {#5} { none }
    \seq_map_inline:Nn \l__hit_font_families_seq
      { \__hit_font_italic_family:nNNN {##1} #1 #2 #3 }
    \__hit_font_emph_family:NNN #1 #2 #3
    \group_end:
    \normalfont
  }
%    \end{macrocode}
%
% 一个字族。宋体改成斜切，别的字族只补 \texttt{AutoFakeSlant}。
%    \begin{macrocode}
\cs_new_protected:Npn \__hit_font_italic_family:nNNN #1#2#3#4
  {
    \tl_if_in:nnF {#1} { / }
      {
        \prop_get:NnN #2 {#1} \l__hit_font_name_tl
        \prop_get:NnNF #3 {#1} \l__hit_font_options_clist
          { \clist_clear:N \l__hit_font_options_clist }
        \bool_lazy_or:nnTF
          { \str_if_eq_p:ee {#1} { \CJKrmdefault } }
          { \str_if_eq_p:nn {#1} { zhsong } }
          { \__hit_font_italic_song: }
          { \__hit_font_italic_other: }
        \tl_if_eq:NNF \l__hit_font_new_options_clist \l__hit_font_options_clist
          {
            \use:x
              {
                \exp_not:N #4 {#1}
                  { \exp_not:V \l__hit_font_new_options_clist }
                  { \exp_not:V \l__hit_font_name_tl }
              }
          }
      }
  }
\cs_new_protected:Npn \__hit_font_strip_italic:NN #1#2
  {
    \clist_clear:N #2
    \clist_map_inline:Nn #1
      {
        \regex_match:nnF
          { ^\s*(ItalicFont|ItalicFeatures|BoldItalicFont|BoldItalicFeatures|AutoFakeSlant|SlantFont|slantfont|Script)\b } {##1}
          { \clist_put_right:Nn #2 {##1} }
      }
  }
\cs_new_protected:Npn \__hit_font_italic_song:
  {
    \__hit_font_strip_italic:NN \l__hit_font_options_clist \l__hit_font_new_options_clist
    \clist_put_right:Nx \l__hit_font_new_options_clist
      { AutoFakeSlant = \c_docxlike_font_fake_slant_tl }
  }
\cs_new_protected:Npn \__hit_font_italic_other:
  {
    \clist_set_eq:NN \l__hit_font_new_options_clist \l__hit_font_options_clist
    \bool_set_false:N \l__hit_font_has_italic_bool
    \clist_map_inline:Nn \l__hit_font_options_clist
      {
        \regex_match:nnT
          { ^\s*(ItalicFont|AutoFakeSlant|SlantFont|slantfont)\b } {##1}
          { \bool_set_true:N \l__hit_font_has_italic_bool }
      }
    \bool_if:NF \l__hit_font_has_italic_bool
      {
        \clist_remove_all:Nn \l__hit_font_new_options_clist { Script = CJK }
        \clist_put_right:Nx \l__hit_font_new_options_clist
          { AutoFakeSlant = \c_docxlike_font_fake_slant_tl }
      }
  }
%    \end{macrocode}
%
% \cs{emph} 用的那一族。
%    \begin{macrocode}
\cs_new_protected:Npn \__hit_font_emph_family:NNN #1#2#3
  {
    \bool_gset_false:N \g__hit_font_emph_bool
    \str_case:VnF \g__hit_font_emph_songti_tl
      {
        { fake-slant } { }
        { kaiti }
          {
            \prop_get:NnNT #1 { zhkai } \l__hit_font_emph_name_tl
              {
                \prop_get:NnNF #2 { zhkai } \l__hit_font_options_clist
                  { \clist_clear:N \l__hit_font_options_clist }
                \__hit_font_emph_declare:N #3
              }
          }
      }
      {
        \tl_set_eq:NN \l__hit_font_emph_name_tl \g__hit_font_emph_songti_tl
        \clist_clear:N \l__hit_font_options_clist
        \__hit_font_emph_declare:N #3
      }
  }
\tl_new:N \l__hit_font_emph_upright_tl
\keys_define:nn {hit/emph-font-face}
  {UprightFont .tl_set:N=\l__hit_font_emph_upright_tl,unknown .code:n={}}
\cs_new_protected:Npn \__hit_font_emph_declare:N #1
  {
    \__hit_font_strip_italic:NN \l__hit_font_options_clist \l__hit_font_emph_options_clist
    % 文件预设的族名不一定可由系统识别；强调使用该族常规字面的文件模式。
    \tl_set_eq:NN\l__hit_font_emph_upright_tl\l__hit_font_emph_name_tl
    \exp_args:NnV\keys_set:nn{hit/emph-font-face}\l__hit_font_emph_options_clist
    \clist_put_right:Nx \l__hit_font_emph_options_clist
      { ItalicFont = \l__hit_font_emph_upright_tl , AutoFakeSlant = 0 }
    \use:x
      {
        \exp_not:N #1 { hitemphsong }
          { \exp_not:V \l__hit_font_emph_options_clist }
          { \exp_not:V \l__hit_font_emph_name_tl }
      }
    \bool_gset_true:N \g__hit_font_emph_bool
  }
%    \end{macrocode}
% \end{macro}
%
% \begin{macro}{\emph}
% \cs{emph} 在 \cs{em} 之后看一眼：字形已经变斜、中文字族是宋体的，中文换到
% \texttt{hitemphsong}；嵌套的 \cs{emph} 把字形变回正体时，中文也回到进来之前的宋体那一族。
% 别的字族不动，\cs{em} 的斜体已经把它们斜切了。当前字族两个引擎各记在
% \cs{l\_xeCJK\_family\_tl} 与 \cs{l\_ctex\_ltj\_family\_tl} 里，\hologo{pdfTeX} 下是 \pkg{CJK} 的 \cs{CJK@family}。\cs{textit} 不经过这里。
%    \begin{macrocode}
\tl_new:N \l__hit_emph_family_tl
\tl_new:N \l__hit_emph_song_tl
\cs_new:Npn \__hit_emph_current_family:
  {
    \cs_if_exist:NTF \l_xeCJK_family_tl
      { \l_xeCJK_family_tl }
      {
        \sys_if_engine_luatex:TF
          { \l__hit_cjk_family_tl }
          { \cs_if_exist:NT \CJK@family { \CJK@family } }
      }
  }
\cs_new_protected:Npn \__hit_emph_cjk:
  {
    \tl_set:Nx \l__hit_emph_family_tl { \__hit_emph_current_family: }
    \str_if_eq:eeTF { \f@shape } { it }
      {
        \bool_lazy_or:nnT
          { \str_if_eq_p:ee { \l__hit_emph_family_tl } { \CJKrmdefault } }
          { \str_if_eq_p:ee { \l__hit_emph_family_tl } { zhsong } }
          {
            \tl_set_eq:NN \l__hit_emph_song_tl \l__hit_emph_family_tl
            \CJKfamily { hitemphsong }
          }
      }
      {
        \str_if_eq:eeT { \l__hit_emph_family_tl } { hitemphsong }
          { \exp_args:NV \CJKfamily \l__hit_emph_song_tl }
      }
  }
\AddToHook { begindocument / end }
  {
    \bool_if:NT \g__hit_font_emph_bool
      {
        \NewCommandCopy \__hit_emph_orig:n \emph
        \RenewDocumentCommand \emph { m } { \__hit_emph_orig:n { \__hit_emph_cjk: #1 } }
      }
  }
%    \end{macrocode}
% \end{macro}
%
% MS~Word 里的 1\,pt 就是 \hologo{LaTeX} 里的 1\,bp，而 1\,bp 等于 \hologo{LaTeX} 里的
% 72.27/72\,pt。所以 Word 上看到的磅数，在这里直接写成 bp 即可。对照如下：
% \begin{longtblr}[caption={字号名与磅数对照}, label={tab:impl-font-size}]
%   {colspec={l l l l}}
% \toprule
% 字号 & Word 的磅 & 毫米 & \hologo{TeX} 的磅 \\ \midrule
% 初号 & 42bp & 14.82mm & 42.1575pt \\
% 小初 & 36bp & 12.70mm & 36.135 pt \\
% 一号 & 26bp & 9.17mm & 26.0975pt \\
% 小一 & 24bp & 8.47mm & 24.09pt \\
% 二号 & 22bp & 7.76mm & 22.0825pt \\
% 小二 & 18bp & 6.35mm & 18.0675pt \\
% 三号 & 16bp & 5.64mm & 16.06pt \\
% 小三 & 15bp & 5.29mm & 15.05625pt \\
% 四号 & 14bp & 4.94mm & 14.0525pt \\
% 小四 & 12bp & 4.23mm & 12.045pt \\
% 五号 & 10.5bp & 3.70mm & 10.59375pt \\
% 小五 & 9bp & 3.18mm & 9.03375pt \\
% 六号 & 7.5bp & 2.56mm & \\
% 小六 & 6.5bp & 2.29mm & \\
% 七号 & 5.5bp & 1.94mm & \\
% 八号 & 5bp & 1.76mm & \\
% \bottomrule
% \end{longtblr}
%
% 定义一个字号命令，参数为名称与磅数。字号与行距解耦，行距由可选参数给出。
% 下面十七个字号的取值，\hithesis{} 与 \CTeX、MS~Word、WPS、Adobe
% 的选择一致，模板自己只用到其中十一个，其余留给用户。
%
% 可选参数按自然单倍行高计算倍数行距；省略参数时保留 1.3 倍字号的默认值。
%
% 类自己的代码里用到 \oarg{n} 的地方（报告封面、内封、声明页的字样、报告页眉）写的
% 都是老意思，所以那些地方不再写方括号，改调 \cs{hit@fontsize@x}\marg{名}\marg{n}：
% 永远是 $n \times$ 字号。方括号那条路只留给用户。
%
% 磅数除了包进命令体，还存一份常量 \cs{c\_\_hit\_font\_size\_}\meta{名}\cs{\_dim}。
% \option{styles} 里 \option{font} 下的 \option{size} 键收 \texttt{sanhao} 这种名字，要把它换算成长度，
% 读的就是这份常量。不另抄一张表，改磅数只改这里一处。
%
% 自然行高的系数住在 \pkg{docxlike} 里，这里在调用时才探，探不到就按 $n \times$ 字号算，
% 类装到一半时字号命令照样能用。
%    \begin{macrocode}
\cs_new_protected:Npn \hit@fontsize@x #1#2
  {
    \fontsize { \dim_use:c { c__hit_font_size_#1_dim } }
      { #2 \use:c { c__hit_font_size_#1_dim } }
    \selectfont
  }
\cs_new_protected:Npn \__hit_font_size_word:nn #1#2
  {
    \fontsize { \dim_use:c { c__hit_font_size_#1_dim } }
      {
        \dim_eval:n
          {
            #2 \dimexpr \c_docxlike_format_natural_tl
              \use:c { c__hit_font_size_#1_dim } \relax
          }
      }
    \selectfont
  }
\prg_new_protected_conditional:Npnn \__hit_font_if_word_layout: { TF }
  {
    \tl_if_exist:NTF \c_docxlike_format_natural_tl
      { \prg_return_true: } { \prg_return_false: }
  }
\cs_new_protected:Npn \__hit_define_font_size:nn #1#2
  {
    \dim_const:cn { c__hit_font_size_#1_dim } {#2}
    \exp_args:Nc \NewDocumentCommand {#1} { o }
      {
        \IfNoValueTF {##1}
          { \fontsize {#2} { 1.3 \dimexpr #2 } \selectfont }
          {
            \__hit_font_if_word_layout:TF
              { \__hit_font_size_word:nn {#1} {##1} }
              { \fontsize {#2} { ##1 \dimexpr #2 } \selectfont }
          }
      }
  }
%    \end{macrocode}
%
% \begin{macro}{\chuhao}
% \begin{macro}{\xiaochu}
% \begin{macro}{\yihao}
% \begin{macro}{\xiaoyi}
% \begin{macro}{\erhao}
% \begin{macro}{\xiaoer}
% \begin{macro}{\sanhao}
% \begin{macro}{\xsanhao}
% \changes{v3.1b}{2022/06/06}{增加 \cs{xsanhao} 为了深圳校区的原创性声明标题}
% \begin{macro}{\xiaosan}
% \begin{macro}{\sihao}
% \begin{macro}{\xiaosi}
% \begin{macro}{\wuhao}
% \begin{macro}{\xiaowu}
% \begin{macro}{\liuhao}
% \begin{macro}{\xiaoliu}
% \begin{macro}{\qihao}
% \begin{macro}{\bahao}
% 一组字号定义。
%    \begin{macrocode}
\__hit_define_font_size:nn{chuhao}{42bp}
\__hit_define_font_size:nn{xiaochu}{36bp}
\__hit_define_font_size:nn{yihao}{26bp}
\__hit_define_font_size:nn{xiaoyi}{24bp}
\__hit_define_font_size:nn{erhao}{22bp}
\__hit_define_font_size:nn{xiaoer}{18bp}
\__hit_define_font_size:nn{sanhao}{16bp}
\__hit_define_font_size:nn{xsanhao}{15.9bp}
\__hit_define_font_size:nn{xiaosan}{15bp}
\__hit_define_font_size:nn{sihao}{14bp}
\__hit_define_font_size:nn{xiaosi}{12bp}
\__hit_define_font_size:nn{wuhao}{10.5bp}
\__hit_define_font_size:nn{xiaowu}{9bp}
\__hit_define_font_size:nn{liuhao}{7.5bp}
\__hit_define_font_size:nn{xiaoliu}{6.5bp}
\__hit_define_font_size:nn{qihao}{5.5bp}
\__hit_define_font_size:nn{bahao}{5bp}
%    \end{macrocode}
%
% \cs{xsanhao} 比三号略小，给深圳校区本科生的原创性声明标题用，取自
% \hitszthesis{}。规范里那个标题是三号，但 \cs{sanhao} 一行排不下。
%
% \subsection{封面上的美术字}
%
% cover-script 选择实际字体或轮廓绘制；auto 优先使用已安装的字体。
%
% 探测走 \cs{fontspec\_font\_if\_exist:nTF}，与上面挑英文字体是同一套。不查
% \cs{lishu} 那种字体命令：那个只在 \option{fontset} 取 \texttt{windows} 时才有定义，
% 而字体装没装上与 \option{fontset} 选哪套是两回事。
%    \begin{macrocode}
%    \end{macrocode}
%
%    \begin{macrocode}
\bool_new:N \g__hit_lishu_bool
%    \end{macrocode}
%
% 显式传入 newCJKfontfamily 的末尾可选参数，防止参数扫描继续展开调用后的代码。
% 隶书和华文新魏缺少粗体字面，需启用合成粗体；Lua 下也须向 NFSS 登记粗体形状。
% 字体探测状态是全局的，创建的字体命令也必须全局保存，否则首次调用所在组结束后命令会失效。
% pdfTeX 无这些字体的度量，直接使用轮廓替代。
%    \begin{macrocode}
\cs_new_protected:Npn \__hit_cover_font:nnN #1#2#3
  {
    \sys_if_engine_pdftex:F
      {
        \str_if_eq:nnF{#1}{outline}
        {\fontspec_font_if_exist:nT {#1}
          {
            \bool_gset_true:N #3
            \tl_set:Nn\l__hit_font_role_name_tl{#1}
            \tl_set:Nn\l__hit_font_role_options_tl{AutoFakeBold=\c_docxlike_font_fake_bold_tl}
            \str_if_eq:nnTF{#2}{hit@cover@lishu}
              {\__hit_font_baseline_options:n{lishu}}
              {\__hit_font_baseline_options:n{xinwei}}
            \use:x
              {
                \exp_not:N \newCJKfontfamily \exp_not:c {#2} { \exp_not:n {#1} }
                  [ \exp_not:V\l__hit_font_role_options_tl ]
              }
            \cs_gset_eq:cc {#2} {#2}
          }}
      }
  }
%    \end{macrocode}
%
% 仅在封面实际请求该字体时探测，并缓存结果。XeTeX 对缺失字体的查找可能触发
% mktextfm，因此不应在每份文档加载时无条件探测。
%    \begin{macrocode}
\bool_new:N \g__hit_lishu_probed_bool
\cs_new_protected:Npn \__hit_lishu_probe:
  {
    \bool_if:NF \g__hit_lishu_probed_bool
      {
        \bool_gset_true:N \g__hit_lishu_probed_bool
        \exp_args:Ne\__hit_cover_font:nnN {\__hit_font_role_or:nn{lishu}{LiSu}} { hit@cover@lishu } \g__hit_lishu_bool
      }
  }
%    \end{macrocode}
%
% \begin{macro}{\xinwei}
% 深圳本科报告的校名与文种使用华文新魏，首次使用时探测并缓存结果。
% 缺少字体时用宋体计算行高、以轮廓绘制文字，见 \file{src/pages/hit-report-cover-shenzhen-bachelor.dtx}。
%    \begin{macrocode}
\bool_new:N \g__hit_xinwei_bool
\bool_new:N \g__hit_xinwei_probed_bool
\cs_new_protected:Npn \__hit_xinwei_probe:
  {
    \bool_if:NF \g__hit_xinwei_probed_bool
      {
        \bool_gset_true:N \g__hit_xinwei_probed_bool
        \exp_args:Ne\__hit_cover_font:nnN {\__hit_font_role_or:nn{xinwei}{STXinwei}} { hit@cover@xinwei } \g__hit_xinwei_bool
      }
  }
\cs_new_protected:Npn \xinwei
  {
    \__hit_xinwei_probe:
    \bool_if:NTF \g__hit_xinwei_bool { \hit@cover@xinwei } { \songti }
  }
%    \end{macrocode}
%
% 待导言区的 ctex 字体集加载完成后，再补缺失的 lishu 命令，避免与
% ctex-fontset-windows.def 的 NewDocumentCommand 定义冲突。
%    \begin{macrocode}
\AddToHook { begindocument / before }
  {
    \cs_if_exist:NF \lishu
      {
        \cs_new_protected:Npn \lishu
          {
            \__hit_lishu_probe:
            \bool_if:NTF \g__hit_lishu_bool { \hit@cover@lishu } { \songti }
          }
      }
  }
%    \end{macrocode}
% \end{macro}
%
% \begin{macro}{\hit@cover@script:nnn}
% \marg{字体}\marg{排字体的代码}\marg{照轮廓画的代码}，字体写 \texttt{lishu} 或
% \texttt{xinwei}。\option{cover-script} 写 \texttt{font} 而机器上没这个字体时给一句提示
% 再画：钉死了却办不到，不说一声排出来会与预期两样。
%    \begin{macrocode}
\cs_new_protected:Npn \hit@cover@script:nnn #1#2#3
  {
    \use:c { __hit_#1_probe: }
    \str_case:onF { \hit@cover@script }
      {
        { outline } {#3}
        { font }
          {
            \bool_if:cTF { g__hit_#1_bool } {#2}
              { \msg_warning:nn { hithesis } { cover-script-no-font } #3 }
          }
      }
      { \bool_if:cTF { g__hit_#1_bool } {#2} {#3} }
  }
%    \end{macrocode}
% \end{macro}
%
% \begin{macro}{\hit@cover@bottommark}
% 落款那处美术字的排法：隶书小二，粗细随段落样式。
%
% 要把 \cs{CJKglue} 归零。\pkg{xeCJK} 在正文里给汉字之间留一段可伸缩的胶
% 好让整行对齐，美术字这里不该有，原件本来也是紧排的。
%
% 字距必须在 selectfont 之后清零，否则会被字号切换重设。
%
% 照轮廓画的那条路每字后面补字号的 1/280：\hologo{LuaLaTeX} 上断行层给合成粗体的字
% 每字加这么多（Word 的规矩，见 \pkg{docxlike} 的“字符格式”），轮廓盒不是字形，
% 断行层不管，这里自己补，两条路才一样宽。\hologo{XeLaTeX} 上字体那条路也不加。
% 词条改成了轮廓里没有的字，就排宋体。
%    \begin{macrocode}
\cs_new_protected:Npn \hit@cover@bottommark
  {
    \group_begin:
      \hit@cover@script:nnn { lishu }
        {
          \hit@cover@lishu
          \fontsize { 18bp } { 18bp } \selectfont
          \hit@engine@cjkglue@zero@cc:
          \hit@term@cover@bottommark
        }
        {
          \fontsize { 18bp } { 18bp } \selectfont
          \hit@calligraphy@if@cover:nnTF { lishu } { \hit@term@cover@bottommark }
            {
              \sys_if_engine_luatex:T
                {
                  \str_if_in:NnT \f@series { b }
                    { \dim_set:Nn \l__hit_calligraphy_extra_dim { 18bp / 280 } }
                }
              \hit@calligraphy:nn { lishu } { \hit@term@cover@bottommark }
            }
            {
              \songti \selectfont
              \hit@engine@cjkglue@zero@cc:
              \hit@term@cover@bottommark
            }
        }
    \group_end:
  }
%    \end{macrocode}
% \end{macro}
%    \begin{macrocode}
%</shared-fonts>
%    \end{macrocode}
%
% \end{macro}
% \end{macro}
% \end{macro}
% \end{macro}
% \end{macro}
% \end{macro}
% \end{macro}
% \end{macro}
% \end{macro}
% \end{macro}
% \end{macro}
% \end{macro}
% \end{macro}
% \end{macro}
% \end{macro}
% \end{macro}
% \end{macro}
%
% \Finale
